\section{Conclusion}

We set out to validate that incremental SMT verification would achieve 2-5x speedup (predicted, not validated) over batch re-verification for LLM code repair in typed Python with Pydantic. Instead, our validation pipeline failed before generating a single line of code, revealing that infrastructure robustness---not just theoretical soundness---determines whether a hypothesis can be tested. This paper presents the hypothesis, the implemented experimental pipeline, and the lessons learned from complete validation failure due to API authentication errors.

The hypothesis---that incremental SMT verification (re-verifying only modified functions + dependencies) reduces verification time by 2-5x (predicted, not validated) compared to batch re-verification---remains \textbf{theoretically plausible but empirically unvalidated}. All experimental validation was blocked at the foundational hypothesis (h-e1: static analyzers can extract SMT constraints from LLM-generated typed code) due to 100\% API authentication failure rate (48/48 generation attempts returned HTTP 401 ``API key is invalid'').

\subsection{The Artifact Contribution}

Despite generation failure, the dataset extension pipeline succeeded: HumanEval was mechanically extended with 100 Pydantic-annotated prompts. This artifact addresses a gap in formal verification research---prior benchmarks (HumanEval, APPS, CodeContests) lack type annotations required for SMT constraint extraction. The template-based transformation approach (function signature + docstring → Pydantic BaseModel + validators) scales mechanically, enabling future research on static analysis for LLM code.

The constraint extraction pipeline (DatasetExtender → LLMGenerator → PyreExtractor → Z3Validator → MetricsEngine) was validated on error files, confirming architectural soundness on negative cases (pipeline processes malformed inputs without crashing; positive validation requires actual LLM code). The modular design allowed partial validation despite generation failure---dataset extension succeeded independently, preventing cascading failures.

\subsection{Infrastructure Lesson}

Pre-flight validation is critical for experiments with external API dependencies. A single test API call before starting batch experiments would have detected the auth error immediately. Experimental design must address failure modes explicitly: pre-flight validation, fast-fail logic for non-retryable errors (authentication, invalid input), and checkpoint/resume for long-running experiments.

\subsection{Honest Assessment of Claims}

We cannot claim:
\begin{itemize}
\item ``Incremental SMT achieves 2-5x speedup'' (prediction P1 untested)
\item ``Speedup scales with codebase size'' (P2 untested)
\item ``Conservative dependency analysis maintains soundness'' (P3 untested, design intent only)
\item ``Type annotations enable constraint extraction from LLM code'' (h-e1 foundational claim untested)
\end{itemize}

All quantitative speedup predictions await validation pending infrastructure fixes (valid API key) and completion of the h-e1 → h-m1 → h-m2 → h-m3 hypothesis chain.

\subsection{Limitations and Boundary Conditions}

\textbf{Single language scope:} Only Python + Pydantic + Pyre was designed for testing. Claims about ``statically-typed languages'' (plural) are overstated. Generalization to Rust + Prusti or other languages untested.

\textbf{Untested assumptions:} Assumption A1 (LLM code has extractable constraints) and A2 (repair locality transfers from human to LLM code) are unverified. If A1 fails, the approach requires LLM fine-tuning or pivot to neural extraction. If A2 fails (LLM repairs are broad, not localized), speedup benefits may disappear even if extraction succeeds.

\textbf{No baseline comparison:} The 2-5x speedup figure is speculative (theoretical analysis), not measured. Even if extraction works, actual speedup could be marginal (1.1-1.3x) rather than predicted (2-5x).

\subsection{Future Work}

\textbf{Immediate:} Retry h-e1 with valid Anthropic API key. If extraction\_rate $\geq$90\%, the hypothesis validation can proceed to mechanism testing (h-m1: constraint extraction pipeline, h-m2: dependency analysis, h-m3: speedup measurement). If extraction\_rate <90\%, pivot to neural constraint extraction.

\textbf{Medium-term:} Extend approach to Rust + Prusti for language generalization. Test whether incremental SMT transfers beyond Python. Validate repair locality assumption (A2) on LLM code---neural repair patterns may differ from human edits documented in CURE/CoCoNut literature.

\textbf{Long-term vision:} Standardize typed benchmarks for formal verification research across multiple languages (Python + Pydantic, Rust + Prusti, TypeScript + type predicates). Establish evaluation protocol for comparing neural code generation with/without SMT verification. Enable practical deployment of formally-verified LLM code repair in safety-critical domains.

\subsection{Closing the Loop}

We began with the question: can infrastructure robustness be separated from theoretical soundness in hypothesis testing? The answer is no---infrastructure is prerequisite, not afterthought. A scientifically sound hypothesis with elegant decomposition (existence → mechanisms), rigorous experimental protocol, and complete implementation can still fail to advance knowledge if external dependencies (API access, network reliability, service availability) are not validated before execution. The HumanEval + Pydantic dataset extension is the artifact from this validation attempt---a resource for future constraint extraction research, and a reminder that experimental methodology extends beyond hypothesis design to encompass infrastructure reliability.
